% =====================================================================
% PRISM Phase 0 -- rewritten capability matrix (drop-in Table 1)
% Rows are ARCHITECTURAL capabilities, each reproducible from the method
% definition. No empirical property (calibration, accuracy) appears here.
% Requires: \usepackage{booktabs,pifont}
%   \newcommand{\cmark}{\ding{51}}
%   \newcommand{\xmark}{\ding{55}}
% =====================================================================

\begin{table}[t]
\centering
\caption{Architectural capability matrix. Each entry follows from the method
definition, independent of experimental quality. \cmark{}~present, \xmark{}~absent
by default. A support transform (row~3) can be added to \emph{any} method; we do
so for all baselines in the constraint comparison, so feasibility is treated as a
shared, testable property rather than unique to PRISM.}
\label{tab:novelty}
\begin{tabular}{lcccccc}
\toprule
Capability & PRISM & cINN & NPE-CNF & VI & FNO/DeepONet & PINN \\
\midrule
Invertible latent transform      & \cmark & \cmark & \cmark & \cmark & \xmark & \xmark \\
Tractable conditional density    & \cmark & \cmark & \cmark & \cmark & \xmark & \xmark \\
Explicit support transform       & \cmark & \xmark & \xmark & \xmark & \xmark & \xmark \\
Auxiliary physical forward map    & \cmark & \xmark & \xmark & \xmark & \cmark & \cmark \\
Amortized single-pass sampling   & \cmark & \cmark & \cmark & \cmark & \xmark & \xmark \\
Continuous-time (ODE) transport  & \cmark & \xmark & \cmark & \xmark & \xmark & \xmark \\
\bottomrule
\end{tabular}
\end{table}
